% !TeX program = pdflatex
\documentclass[11pt]{article}

\usepackage[margin=1in]{geometry}
\usepackage{lmodern}
\usepackage{microtype}
\usepackage{setspace}
\setstretch{1.08}
\usepackage{amsmath,amssymb,mathtools,bm,booktabs,tabularx,array,amsthm}
\usepackage{enumitem}
\usepackage[round,authoryear]{natbib}
\usepackage{hyperref}
\hypersetup{
  colorlinks=true,
  linkcolor=black,
  citecolor=black,
  urlcolor=black,
  pdftitle={Supplementary Material: Mean-Tilted Intervals: Regression and Dynamic Models},
  pdfauthor={Antonio De Leon, Raquel Prado, and Bruno Sanso}
}

\newcommand{\E}{\mathbb{E}}
\newcommand{\Normal}{\mathcal{N}}
\newcommand{\Exp}{\mathrm{Exp}}
\newcommand{\GIG}{\mathrm{GIG}}
\newcommand{\ind}{\mathbf{1}}
\newcommand{\diag}{\mathrm{diag}}
\newcommand{\vect}[1]{\bm{#1}}
\newcommand{\mat}[1]{\bm{#1}}
\DeclareMathOperator{\argmin}{arg\,min}
\newtheorem{proposition}{Proposition}
\renewcommand{\theproposition}{S\arabic{proposition}}
\renewcommand{\theequation}{S\arabic{equation}}
\renewcommand{\thesection}{S\arabic{section}}
\renewcommand{\thetable}{S\arabic{table}}

\title{Supplementary Material for \emph{Mean-Tilted Intervals:
Regression and Dynamic Models}}
\author{Antonio De Leon \and Raquel Prado \and Bruno Sans\'o}
\date{}

\begin{document}
\maketitle

\begin{abstract}
\noindent
This supplement records computational identities and implementation conditions
for generalized-Bayes endpoint models based on the mean-preserving and
mean-tilted interval losses. It gives the pseudo-asymmetric-Laplace
augmentation, Gaussian root updates, fixed-tilt shifts, ECM mode equations,
root-label summaries, fixed-feature endpoint readouts, dynamic endpoint-state
updates, and diagnostic requirements.
\end{abstract}

\section{Notation and Targets}
\label{sec:supp-targets}

For an observation \(y_i\) and raw roots \(\eta_{1i}\) and \(\eta_{2i}\),
define
\[
e_i=(y_i-\eta_{1i})(y_i-\eta_{2i}),\qquad
\rho_c(r)=r\{c-\ind(r<0)\}.
\]
The MPI loss is \(\rho_c(e_i)\). With fixed tilt \(\delta_i\), the MTI loss is
\[
\rho_c(e_i)-c\delta_i(\eta_{1i}+\eta_{2i}-2y_i).
\]
The cumulative loss over observed indices \(I_{\mathrm{obs}}\) is denoted by
\(L_{c,\delta}(\theta)\). The fixed-target generalized posterior is
\[
\pi(\theta\mid D)
\propto
\exp\{-\omega_{\mathrm R}L_{c,\delta}(\theta)\}\pi_0(\theta).
\]
The symbols \(c\), \(\delta\), and \(\omega_{\mathrm R}\) have different roles:
content, retained-mean tilt, and learning rate, respectively.

\section{Root Labels and Endpoint Summaries}
\label{sec:supp-labels}

The product-residual loss is invariant to exchanging the two raw roots. With
exchangeable root priors, the unordered pair is identified but the labels
``root 1'' and ``root 2'' are not. Endpoint summaries are therefore computed
after sorting the roots at the evaluation points:
\[
L_i=\min(\eta_{1i},\eta_{2i}),\qquad
U_i=\max(\eta_{1i},\eta_{2i}),\qquad
W_i=U_i-L_i.
\]

For static linear roots, define
\[
\beta_m=\frac{\beta_1+\beta_2}{2},\qquad
\beta_s=\frac{\beta_2-\beta_1}{2}.
\]
Then \(M(x)=x^\top\beta_m\), \(S(x)=x^\top\beta_s\), and
\[
L(x)=M(x)-|S(x)|,\qquad
U(x)=M(x)+|S(x)|,\qquad
W(x)=2|S(x)|.
\]
The midpoint coefficient \(\beta_m\) is root-label invariant. The signed gap
coefficient \(\beta_s\) changes sign under a root swap. Lower- and
upper-endpoint coefficient summaries require a declared design region where the
two roots remain ordered.

\section{Pseudo-AL Augmentation}
\label{sec:supp-pseudo-al}

Let
\[
\sigma_{\mathrm R}=\omega_{\mathrm R}^{-1},\qquad
\xi_c=\frac{1-2c}{c(1-c)},\qquad
\phi_c=\frac{2}{c(1-c)}.
\]
The single-observation loss factor satisfies
\begin{align}
&\frac{c(1-c)}{\sigma_{\mathrm R}}
\exp\{-\rho_c(e_i)/\sigma_{\mathrm R}\}\notag\\
&\quad=\int_0^\infty
\Normal(e_i\mid \xi_cv_i,\phi_c\sigma_{\mathrm R}v_i)
\Exp(v_i\mid\mathrm{rate}=1/\sigma_{\mathrm R})\,dv_i .
\end{align}
Thus
\[
v_i\mid\cdot\sim
\GIG\left\{
\frac12,\
\frac{1}{2\sigma_{\mathrm R}c(1-c)},\
\frac{c(1-c)e_i^2}{2\sigma_{\mathrm R}}
\right\}.
\]
The latent scales are part of the computational representation of a
loss-based update. They are not latent response variables.

\section{Gaussian Root Updates and Fixed Tilt}
\label{sec:supp-static-updates}

Let \(\mat X_{\mathrm{obs}}\) be the observed design matrix and let
\(\mat W_v=\diag\{(\phi_c\sigma_{\mathrm R}v_i)^{-1}:i\in I_{\mathrm{obs}}\}\).
Conditional on the second root, define
\[
\mat A_1=\diag(\vect y_{\mathrm{obs}}-\vect\eta_{2,\mathrm{obs}})
\mat X_{\mathrm{obs}},
\qquad
\vect z_1=\vect y_{\mathrm{obs}}\circ
(\vect y_{\mathrm{obs}}-\vect\eta_{2,\mathrm{obs}})-\xi_c\vect v.
\]
For Gaussian prior canonical parameters \((\mat Q_{01},\vect h_{01})\),
\[
\mat Q_1=\mat Q_{01}+\mat A_1^\top\mat W_v\mat A_1,\qquad
\vect h_1=\vect h_{01}+\mat A_1^\top\mat W_v\vect z_1.
\]
Hence
\[
\beta_1\mid\cdot\sim
\Normal(\mat Q_1^{-1}\vect h_1,\mat Q_1^{-1}).
\]
The second root is updated by exchanging labels.

For fixed nonzero tilt, let \(\vect\delta_{\mathrm{obs}}\) be fixed over the
observed design. The precision matrix is unchanged, and the canonical vector is
shifted to
\[
\vect h_1^{\mathrm{MTI}}
=\vect h_1+
\omega_{\mathrm R}c\mat X_{\mathrm{obs}}^\top\vect\delta_{\mathrm{obs}},
\]
with the same shift for the second root. Proper Gaussian priors give a finite
normalizing constant in finite dimension because the tilt is linear in the
roots.

\section{ECM Mode Calculation}
\label{sec:supp-ecm}

The ECM calculation maximizes the fixed generalized-posterior kernel. It uses
the exact inverse-scale moment
\[
\tau_i=E(v_i^{-1}\mid\cdot)=\frac{1}{c(1-c)|e_i|}
\]
away from \(e_i=0\), with a prespecified numerical floor near zero. Holding
\(\beta_2\) fixed, set
\[
\mat B_1=\diag(\vect y_{\mathrm{obs}}-\vect\eta_{2,\mathrm{obs}})
\mat X_{\mathrm{obs}},
\qquad
\vect r_1=\vect y_{\mathrm{obs}}^{\,2}
-\vect y_{\mathrm{obs}}\circ\vect\eta_{2,\mathrm{obs}}.
\]
With ridge prior precision \(\mat P_1\), the conditional maximization solves
\[
\left\{\mat P_1+
\frac{1}{\phi_c\sigma_{\mathrm R}}\mat B_1^\top\diag(\vect\tau)\mat B_1
\right\}\beta_1
=
\frac{1}{\phi_c\sigma_{\mathrm R}}
\mat B_1^\top\{\diag(\vect\tau)\vect r_1-\xi_c\vect 1\}
+\omega_{\mathrm R}c\mat X_{\mathrm{obs}}^\top\vect\delta_{\mathrm{obs}}.
\]
The second root uses the same equation with labels exchanged. A full ECM cycle
updates the latent-scale moments, solves the two conditional Gaussian systems,
and accepts the result only if the exact penalized objective decreases.

\section{Conditionally Gaussian Shrinkage}
\label{sec:supp-shrinkage}

The zero-tilt MPI sampler can use a conditionally Gaussian regularized
horseshoe prior through the Nishimura--Suchard representation
\citep{CarvalhoPolsonScott2010HS,PiironenVehtari2017RHS,NishimuraSuchard2023SSS}.
Conditional on the local, global, and slab-scale states, each root coefficient
block has a Gaussian prior precision and therefore fits directly into the
Gaussian update above. This construction is stated for MPI. Extending nonzero
tilt through random shrinkage scales requires its own propriety and diagnostic
checks.

\section{Deterministic Feature Readouts}
\label{sec:supp-features}

Let \(\mat X_{\mathrm F}\) denote a deterministic feature matrix obtained before
the endpoint update is fitted. The static formulas apply with
\(\mat X_{\mathrm{obs}}\) replaced by the observed rows of
\(\mat X_{\mathrm F}\). Echo-state and deep echo-state features
\citep{Jaeger2001EchoState,GallicchioMicheliPedrelli2018DeepESNDesign} are one
source of such a matrix. The recurrent representation is part of the design
construction; posterior uncertainty in the endpoint model is conditional on the
chosen features unless the feature-generation mechanism is included in the
model.

\section{Dynamic Endpoint-State Updates}
\label{sec:supp-dynamic}

A dynamic endpoint model replaces static coefficients by root-specific state
trajectories,
\[
\eta_{jt}=f_t^\top\theta_{jt},\qquad
\theta_{jt}=G_t\theta_{j,t-1}+\varepsilon_{jt},\qquad j=1,2.
\]
Gaussian initial and evolution distributions give a Gaussian prior over each
root path. Because the product residual contains both endpoint paths, the joint
augmented log density is quartic in the two paths. Conditional on one complete
path, however, the other path has a linear Gaussian observation equation.
Root-specific forward-filtering backward-sampling can therefore be alternated
within the specified fixed-target generalized posterior
\citep{GoncalvesMigonBastos2020DQLM}.

Discount factors and component scales are part of the state prior. Fixed values
or precomputed templates define a fixed target. Adaptive choices are best
reported as sensitivity analyses unless they are embedded in a fully specified
probability model.

\section{Diagnostics and Evidence Boundaries}
\label{sec:supp-diagnostics}

MCMC summaries should report multiple-chain convergence diagnostics, effective
sample sizes, Monte Carlo error, root-label handling, and endpoint functional
summaries. ECM summaries should report objective descent, iteration counts,
stopping criteria, final roots, fitted content, tilt, and sensitivity to
starting values. Fixed-feature analyses should report how the features were
chosen. Dynamic analyses should report state-prior choices, smoothing
diagnostics, and whether discount or scale values were fixed in advance.

These diagnostics assess computation for a fixed endpoint target. They do not
establish tolerance validity or response-predictive calibration. Such claims
require a separate validation design aligned with the intended statistical
statement.

\bibliographystyle{plainnat}
\bibliography{refs}

\end{document}
